\chapter{Couche web : contrôleurs MVC2}\label{chap:web}

\section{Le patron MVC2 appliqué}
Dans MVC2, le contrôleur est une servlet. Elle reçoit la requête, appelle le modèle (les services), dépose le résultat
dans la requête, puis transmet à une JSP qui se contente d'afficher. Les JSP sont placées sous \code{WEB-INF/views} :
le navigateur ne peut donc jamais les demander directement, seule une servlet peut les afficher.

J'ai retenu \textbf{une servlet par domaine de gestion}, avec un paramètre \code{action} qui distingue les opérations
(tableau~\ref{tab:servlets}). Une servlet par action multiplierait les classes ; un contrôleur unique deviendrait trop long.

\begin{table}[htbp]
\centering
\caption{Contrôleurs, adresses et actions}\label{tab:servlets}
\begin{tabularx}{\linewidth}{@{}llXl@{}}
\toprule
\textbf{Servlet} & \textbf{Adresse} & \textbf{Actions} & \textbf{Accès}\\
\midrule
\code{VitrineServlet} & \code{/vitrine} & GET : \code{liste}, \code{detail} & public\\
\code{ClientServlet} & \code{/client} & GET : \code{connexion}, \code{inscription} ; POST : \code{connexion}, \code{inscription}, \code{deconnexion} & public\\
\code{PanierServlet} & \code{/panier} & GET : \code{voir}, \code{resume} ; POST : \code{ajouter}, \code{modifier}, \code{retirer}, \code{vider} & client\\
\code{ProduitServlet} & \code{/admin/produits} & GET : \code{liste}, \code{nouveau}, \code{editer} ; POST : \code{enregistrer}, \code{basculer}, \code{supprimer} & admin\\
\code{CategorieServlet} & \code{/admin/categories} & GET : \code{liste}, \code{nouveau}, \code{editer} ; POST : \code{enregistrer}, \code{supprimer} & admin\\
\bottomrule
\end{tabularx}
\end{table}

Une règle simple s'applique partout : \textbf{une action qui modifie des données passe par POST}, jamais par GET. Un lien
ou une image placés sur un autre site ne peuvent ainsi pas déclencher une suppression.

\section{Classe de base}
Les cinq servlets héritent de \code{BaseServlet}, qui factorise ce qui se répète : lecture de l'action, lecture
sécurisée des identifiants et des entiers, affichage d'une vue, redirection, réponse JSON et reconnaissance d'une
requête AJAX (en-tête \code{X-Requested-With}). Une erreur de saisie (identifiant non numérique, par exemple) y est
traitée une fois, au lieu d'être dispersée.

\section{Redirection après POST et messages éphémères}
Après un POST réussi, l'application ne renvoie pas directement une page : elle \textbf{redirige} vers une page GET.
C'est le motif Post/Redirect/Get. Sans lui, un rafraîchissement du navigateur renverrait le formulaire, et ajouterait
une seconde fois le produit au panier.

La redirection fait perdre la requête, donc le message de confirmation (« Produit ajouté ») doit survivre un aller-retour.
La classe \code{Flash} le range en session, et un filtre le transfère à la requête suivante, puis l'efface : le
message ne s'affiche qu'une fois.

\section{Filtres}
Deux filtres s'exécutent avant les servlets, dans cet ordre.

\begin{enumerate}
  \item \code{SecurityFilter} sur toutes les adresses : encodage UTF-8, en-têtes de sécurité, contrôle du jeton CSRF
        sur les requêtes POST, transfert du message éphémère.
  \item \code{AuthFilter} sur \code{/panier} et \code{/admin/*} : exige un client connecté pour le panier, un administrateur
        pour l'administration.
\end{enumerate}

L'accès est contrôlé \textbf{par adresse, dans un filtre}, et non dans chaque servlet. Une nouvelle page d'administration est
protégée dès qu'elle est placée sous \code{/admin}, sans qu'il faille y penser. Un visiteur non connecté est redirigé vers la
connexion, avec mémorisation de la page demandée ; un client qui tente d'ouvrir l'administration reçoit une erreur 403.

\section{Réponses AJAX}
Le mini-panier appelle \code{PanierServlet} en AJAX. La servlet reconnaît la requête et répond en JSON ; sans AJAX, elle redirige
comme n'importe quelle page. Le JSON est écrit à la main, sans bibliothèque, car le contenu est petit : état du
panier, nombre d'articles, total déjà formaté, lignes. Les montants arrivent sous forme de texte formaté : le navigateur n'a
\textbf{aucun calcul d'argent} à faire. Les caractères \code{<}, \code{>} et \code{\&} sont échappés dans le JSON, au cas
où il serait un jour inséré dans du HTML.

Pour un visiteur non connecté, la réponse est un JSON \code{connexion requise} avec l'adresse de la page de connexion,
que le script suit.
